\section{Exact obstruction data}\label{app:data}

These data implement the cusp conditions of \cref{sec:cusp}. The larger polynomial expressions are recorded in the accompanying data (\cref{app:verification}).

\subsection{Rank three}\label{app:r3}
For $A=T_3$, write $s=\sqrt5$. Direct substitution gives
\[
Q=\Bigl(\frac{3-s}2,\ s-2,\ \frac{3-s}4\Bigr),\qquad \xi=\Bigl(\frac{s-1}2,\ \frac{s-1}4,\ s-2\Bigr),
\]
and
\[
M=\begin{pmatrix}-\frac32+\frac{9s}{10}&-1+\frac{3s}5&\frac12-\frac s{10}\\[2pt] -1+\frac{3s}5&\frac{2s}5&\frac12+\frac s{10}\\[2pt] \frac12-\frac s{10}&\frac12+\frac s{10}&\frac12+\frac{2s}5\end{pmatrix}.
\]
These saddle-point values and the first-order condition already occur in \cite[Section~4.1]{MW24}. The recurrence of \cref{sec:cusp} is used with $N=3$, so the coefficients lie in $\Q(\sqrt5)[b_1,b_2,b_3]$. The first-order calculation yields
\begin{equation}\label{eq:r3first}
C=C_0(B)+\frac{\sqrt5}{20}P(B),
\end{equation}
where
\begin{align*}
80C_0(B)&=-60b_1^2-80b_1b_2+40b_1b_3-72b_1+40b_2b_3-68b_2+20b_3^2+8b_3-7,\\
P(B)&=9b_1^2+12b_1b_2-2b_1b_3+8b_1+4b_2^2+2b_2b_3+7b_2+4b_3^2-2b_3.
\end{align*}
Since $B$ and $C$ are rational, $P(B)=0$ and $C=C_0(B)$. For rational $B$, both the rational and the $\sqrt5$ components of $\kappa_2$ and $\kappa_3$ must also vanish; these are four further polynomial equations.

After applying \eqref{eq:r3coords} and clearing denominators, the five equations are $E_0=\cdots=E_4=0$, where
\begin{align*}
E_0&=9u^2-6uv+10uz+8u+v^2-v,\\
E_1&=7200u^3-6000u^2v+8400u^2z+16728u^2+1200uv^2-3600uvz-5832uv\\
&\qquad+2400uz^2+8640uz+9336u-48v^2+48v-100z^3+25z,\\
E_2&=-812u^3+672u^2v-960u^2z-1866u^2-144uv^2+420uvz+654uv\\
&\qquad-300uz^2-960uz-1039u-4v^3+6v^2-2v.
\end{align*}
The two quartics $E_3,E_4$, and all certificate multipliers, are recorded as exact rational expressions in the accompanying data. The precise normalization is
\[
E_0=P(B),\qquad \kappa_2=\frac{E_1}{2400}+\frac{\sqrt5\,E_2}{600},\qquad \kappa_3=\frac{E_3}{2000}+\frac{\sqrt5\,E_4}{6000},
\]
where the $E_i$ are evaluated using \eqref{eq:r3coords}.

\begin{lemma}[Exact elimination]\label{lem:r3elim}
In $\Q[u,v,z]$, the ideal generated by $E_0,\dots,E_4$ equals $\bigl(u,\ v(v-1),\ z(4z^2-1)\bigr)$.
\end{lemma}

\begin{proof}
The certificate gives rational polynomials $A_0,\dots,A_4$ satisfying
\begin{equation}\label{eq:r3cert}
\sum_{i=0}^4A_iE_i=u,
\end{equation}
each product having total degree at most six. Expanding the left-hand side with exact rational coefficients proves \eqref{eq:r3cert}. Consequently $u$ belongs to the obstruction ideal. Next,
\[
E_0\big|_{u=0}=v(v-1),\qquad E_1\big|_{u=0}=-48v(v-1)-25z(4z^2-1),
\]
so the other two displayed generators belong to the ideal as well. Conversely, each $E_i$ reduces to zero under $u=0$, $v^2=v$ and $z^3=z/4$, as verified by direct polynomial division. This proves the reverse containment.
\end{proof}

\subsection{Rank four}\label{app:r4}
Write $s=\sqrt3$. The unique saddle in $(0,1)^4$ is
\[
Q=\Bigl(\frac{s-1}2,\ \frac{3-s}6,\ \frac{2s-3}3,\ \frac{2-s}2\Bigr),
\]
as direct substitution into the four saddle equations confirms. The recurrence of \cref{sec:cusp}, with $r=N=4$, computes $c_1,c_2,c_3,c_4$ over $\Q(\sqrt3)$. The additional logarithmic coefficient is
\[
\kappa_4=c_4-c_1c_3-\tfrac12c_2^2+c_1^2c_2-\tfrac14c_1^4,
\]
and modularity requires $\kappa_4=0$ as well as \eqref{eq:cuspcond}. The first-order condition takes the form $C=C_4(B)+\sqrt3R_4(B)$, where $C_4,R_4\in\Q[b_1,b_2,b_3,b_4]$ and
\begin{align*}
72C_4(B)&=24b_1^2+18b_1b_2+12b_1b_3+6b_1b_4-6b_1+20b_2^2+28b_2b_3+18b_2b_4-4b_2\\
&\qquad+32b_3^2+24b_3b_4-4b_3+24b_4^2+6b_4-9,\\
72R_4(B)&=-3b_1^2+4b_1b_2+4b_1b_3+6b_1b_4+3b_1+7b_2^2+12b_2b_3+14b_2b_4-b_2\\
&\qquad+8b_3^2+32b_3b_4-4b_3+24b_4^2-12b_4.
\end{align*}
Since $B$ and $C$ are rational, $R_4(B)=0$ and $C=C_4(B)$. The rational and $\sqrt3$ components of each of $\kappa_2,\kappa_3,\kappa_4$ also vanish. Let $E^{(4)}_0,\dots,E^{(4)}_6$ be the primitive integer polynomials obtained from these seven rational polynomials in the order $R_4$, then the rational and irrational components for the indices $2,3,4$. Here \emph{primitive} means that denominators are cleared and the common integer factor is removed. Their degrees are $2$, $3$, $3$, $4$, $4$, $5$, $5$, and their exact expressions are recorded in the accompanying data.

Let $J_4=(E^{(4)}_0,\dots,E^{(4)}_6)\subset\Q[b_1,b_2,b_3,b_4]$, and define
\begin{align*}
G_1&=b_1^2-b_1-2b_2b_4+b_2-b_3^2+2b_3+2b_4^2-b_4,\\
G_2&=b_1b_2+b_1+2b_2b_4-b_2+b_3^2-2b_3-2b_4^2+b_4,\\
G_3&=2b_1b_3-2b_1-2b_2b_4-b_2-2b_3^2+b_3+2b_4^2-b_4,\\
G_4&=b_4(b_1+b_2),\\
G_5&=b_2^2-2b_2b_4+2b_2-b_3^2+2b_3+2b_4^2-b_4,\\
G_6&=2b_2b_3+2b_2b_4-b_2+2b_3^2-b_3-2b_4^2+b_4,\\
G_7&=4b_2b_4^2-b_2+4b_3^2-5b_3+2b_4^2-b_4,\\
G_8&=b_3(2b_3-1)(b_3-1),\\
G_9&=b_4(2b_3+2b_4-1),\\
G_{10}&=b_4(2b_4-1)(2b_4+1).
\end{align*}

\begin{lemma}[Rank-four elimination]\label{lem:r4elim}
The polynomials $G_1,\dots,G_{10}$ form a Gr\"obner basis of $J_4$ for the lexicographic order $b_1>b_2>b_3>b_4$. The ideal is radical, and its complex zero set consists of the eight rational vectors of \cref{thm:r4}.
\end{lemma}

\begin{proof}
The accompanying certificate gives rational polynomials $A_{ji}$ such that
\[
G_j=\sum_{i=0}^6A_{ji}E^{(4)}_i\qquad(1\le j\le10),
\]
each product $A_{ji}E^{(4)}_i$ having total degree at most nine. These ten identities hold by exact polynomial multiplication. Conversely, division by the $G_j$ reduces each $E^{(4)}_i$ to zero. All $\binom{10}2=45$ Buchberger $S$-polynomials reduce to zero, which proves the Gr\"obner-basis assertion.

The monomials not divisible by the leading monomials of this basis are
\[
1,\ b_4,\ b_4^2,\ b_3,\ b_3^2,\ b_2,\ b_2b_4,\ b_1.
\]
They form a basis of the quotient ring, which therefore has dimension eight over $\Q$. Direct substitution gives eight distinct rational zeros, namely the vectors of the theorem. Hence there are no further complex zeros and no nontrivial multiplicities, and $J_4$ is radical.
\end{proof}

\subsection{Rank five}\label{app:r5}
Let $t$ be the unique root in $(9/100,1/10)$ of
\[
8t^3-20t^2+12t-1=0;
\]
equivalently, $t=1-\cos(\pi/7)$. The polynomial is irreducible over $\Q$, and direct substitution gives the saddle
\[
Q=\bigl(-4t^2+4t,\ 2t,\ 12t^2-20t+2,\ -8t^2+12t-1,\ t\bigr),
\]
each coordinate of which lies in $(0,1)$. The first three expansion coefficients are computed over $K=\Q(t)$ by the recurrence of \cref{sec:cusp}. Write
\[
c_1-\gamma_0=C_0+tC_1+t^2C_2,\qquad \kappa_j=K_{j,0}+tK_{j,1}+t^2K_{j,2}\quad(j=2,3),
\]
where all component polynomials have rational coefficients in $B$. For rational $B$ and $C$, modularity requires
\[
C_1=C_2=K_{2,0}=K_{2,1}=K_{2,2}=K_{3,0}=K_{3,1}=K_{3,2}=0,\qquad C=C_0.
\]
Clear denominators and remove common integer factors from these eight polynomials, in the displayed order. After the change of coordinates of \cref{thm:r5}, denote them by $E^{(5)}_0,\dots,E^{(5)}_7$. Their degrees are $2$, $2$, $3$, $3$, $3$, $4$, $4$, $4$. Their exact expressions are recorded in the accompanying data.

\begin{lemma}[Rank-five elimination]\label{lem:r5elim}
In $\Q[u,v,x,y,z]$, the obstruction ideal is
\[
\bigl(E^{(5)}_0,\dots,E^{(5)}_7\bigr)=\bigl(u,\ v,\ x^2-x,\ xy-x,\ y^2-y,\ 4z^3-z\bigr).
\]
It is radical and has exactly the nine rational zeros described in \eqref{eq:r5set}.
\end{lemma}

\begin{proof}
For each of the six generators $G_j$ on the right, the certificate records eight rational polynomials $A_{ji}$ such that
\[
G_j=\sum_{i=0}^7A_{ji}E^{(5)}_i,\qquad \deg\bigl(A_{ji}E^{(5)}_i\bigr)\le7.
\]
These identities hold by exact expansion. Conversely, every $E^{(5)}_i$ reduces to zero modulo the six displayed generators. These generators form a Gr\"obner basis for the lexicographic order $u>v>x>y>z$; all fifteen $S$-polynomials reduce to zero. The standard monomials are $\{1,x,y\}\cdot\{1,z,z^2\}$, so the quotient has dimension nine. The three choices of $(x,y)$ in \eqref{eq:r5set}, together with the three distinct roots of $4z^3-z$, give nine rational zeros. This proves radicality and excludes any further complex zeros.
\end{proof}
